\section{Implementación del Subsistema de Visión Artificial}

Este apartado describe la implementación del \textit{pipeline} de visión artificial,
materializando el diseño establecido en la Sección \ref{fig:pipeline_vision}. A
diferencia de soluciones basadas en servicios de terceros o librerías cerradas, los
módulos de detección y reconocimiento óptico de caracteres fueron adaptados y
entrenados a medida para operar con el formato alfanumérico y las condiciones de
captura del proyecto.

\subsection{Adquisición del flujo de video}

La captura de video se implementa como un proceso independiente por cámara
(\texttt{rtsp-capture/capture.py}). El módulo establece una conexión con el flujo
RTSP del grabador digital, aplica un muestreo temporal configurable
(\texttt{DETECTION\_INTERVAL\_MS = 500}) para evitar saturar el bus de datos y
reencamisa las imágenes hacia la cola de procesamiento.

Para tolerar cortes de red o caídas temporales de energía en el recinto, el
capturador incorpora un mecanismo de reconexión automática que reintenta la lectura
sin abortar la ejecución del servicio.

\subsection{Detección de vehículos}

La localización de los automóviles se ejecuta mediante el modelo YOLO26x en su
variante de mayor capacidad. El modelo se carga en memoria empleando el formato de
pesos nativo de PyTorch (\texttt{yolo26x.pt}) o su motor optimizado en TensorRT
(\texttt{yolo26x.engine}) de acuerdo a la disponibilidad de hardware.

La detección se restringe únicamente a las clases de transporte relevantes
(\texttt{vehicle\_classes: [2, 3, 5, 7]}, correspondientes a automóvil, motocicleta,
autobús y camión) con una confianza mínima de $0{,}35$ y un filtro de superficie
que ignora detecciones menores a $3000$ píxeles cuadrados.

\subsection{Seguimiento vehicular}

El seguimiento multiobjeto se implementa integrando el algoritmo ByteTrack
(\texttt{bytetrack.yaml}). El rastreador mantiene la continuidad espacial y
temporal de los vehículos asignando un identificador numérico único (\texttt{track\_id})
que se conserva a través de los fotogramas.

A diferencia del algoritmo SORT convencional, ByteTrack aprovecha tanto las
detecciones de alta confianza como las de baja confianza para recuperar trayectorias
durante oclusiones parciales o momentos en que el vehículo pasa bajo sombras
pronunciadas. El sistema requiere un mínimo de $3$ observaciones válidas
(\texttt{min\_track\_hits}) para confirmar la validez de una trayectoria y acepta una
tolerancia de hasta $30$ fotogramas perdidos (\texttt{grace\_frames}) antes de
eliminar el identificador de la memoria.

\subsection{Detección de placas de matrícula}
\label{subsec:deteccion_placas}

El reconocedor CRNN de la Subsección~\ref{subsec:ocr_crnn} opera sobre un recorte
ajustado de la placa. El trabajador de visión, en cambio, dispone únicamente del
recorte del vehículo que entrega el detector de la Subsección~3.4.2. Entre ambos
media una etapa de localización que el diseño original delegaba en el motor
PaddleOCR; medirla mostró que esa delegación era el eslabón más débil del
\emph{pipeline}, por lo que se entrenó un detector propio.

\subsubsection{Conjunto de datos y revisión humana}

Se partió de los $4.006$ recortes de vehículo del pool congelado. Sobre ellos,
PaddleOCR actuó como pseudo-etiquetador en el sentido estricto del
término: propuso una caja delimitadora por recorte, sin que ninguna de esas cajas
ingresara al conjunto de entrenamiento sin revisión humana.

\begin{itemize}
    \item \textbf{Propuestas automáticas:} $727$ de $4.006$ recortes ($18{,}1\%$)
        contenían una región compatible con una placa. El $81{,}9\%$ restante
        corresponde mayoritariamente a vistas laterales del vehículo, donde la
        placa no es visible; por lo tanto no son fallos del etiquetador sino
        negativos reales.
    \item \textbf{Revisión humana:} se revisaron $1.127$ recortes (todas las
        propuestas más una muestra aleatoria de $400$ negativos). El operador
        aceptó, corrigió o descartó cada caja. El resultado fueron $868$ recortes
        con placa ($77\%$) y $259$ sin placa ($23\%$). De las $700$ propuestas
        automáticas que sí correspondían a una placa, solo $183$ ($26{,}1\%$)
        resultaron geométricamente utilizables sin corrección: el etiquetador
        automático encontraba la placa, pero su caja la recortaba.
    \item \textbf{Participación:} $768$ recortes de entrenamiento, $181$ de
        validación y $178$ de prueba. La partición se realizó por
        \emph{bloque de grabación} y no por trayectoria, de modo que el conjunto
        de prueba contenga horas de captura no vistas durante el entrenamiento.
        No existe ningún bloque compartido entre entrenamiento y evaluación.
\end{itemize}

\subsubsection{Modelo y resultados}

Se ajustó un detector YOLO26n preentrenado en COCO a una única clase
(\texttt{plate}), durante $100$ épocas, con resolución de $640$ píxeles y lote de
$8$. Sobre el conjunto de prueba ($178$ recortes, $140$ placas, $58$ bloques de
grabación no vistos) se obtuvo precisión $0{,}986$, exhaustividad $0{,}971$,
mAP@50 $0{,}990$ y mAP@50--95 $0{,}633$.

La comparación con el etiquetador automático sobre el mismo conjunto de
referencia humana es el argumento que justifica el detector propio: PaddleOCR
alcanza una exhaustividad de $80{,}6\%$ y una precisión de $96{,}3\%$ como
localizador, es decir, no encuentra una de cada cinco placas visibles para el
operador. El detector entrenado eleva la exhaustividad a $97{,}1\%$.

\subsection{Reconocimiento óptico de caracteres (CRNN + CTC)}
\label{subsec:ocr_crnn}

El módulo de reconocimiento de matrículas fue implementado desde cero mediante una
arquitectura profunda combinada \textit{ResNet18 + BiLSTM + CTC}
(\texttt{train\_ocr\_crnn.py}), reemplazando los motores genéricos de terceros por
un modelo adaptado a la tipografía y formato de Bolivia.

\subsubsection{Arquitectura del modelo ResNet-CRNN}

La red se compone de tres bloques funcionales acoplados:

\begin{enumerate}
    \item \textbf{Extractor de características (Backbone CNN):} Basado en ResNet-18
        modificada. Para evitar la pérdida de resolución horizontal crítica en
        secuencias de texto, se alteró el avance (\textit{stride}) de las capas
        convolucionales profundas a $(2, 1)$, reduciendo la altura a $1$ píxel
        mientras se preserva el ancho espacial de $48$ vectores de características.
        Se añade una capa de Normalización por Lotes (\textit{BatchNorm1d}) para
        estabilizar las activaciones antes del bloque recurrente.

    \item \textbf{Modelador de secuencias (RNN):} Red LSTM bidireccional de 2 capas
        (\textit{BiLSTM}) con $512$ unidades ocultas, encargada de modelar las
        dependencias contextuales entre caracteres contiguos de izquierda a derecha.

    \item \textbf{Clasificador y decodificador CTC:} Capa lineal final con una tasa
        de desacople (\textit{Dropout}) de $0{,}3$ conectada a una función de pérdida
        CTC (\textit{Connectionist Temporal Classification}). El alfabeto se define
        con $37$ símbolos: $10$ dígitos ($0$--$9$), $26$ letras mayúsculas ($A$--$Z$)
        y $1$ token de espacio en blanco (\textit{blank}) en el índice $0$.
\end{enumerate}

\subsubsection{Preprocesamiento: normalización de contraste}

Antes de ingresar a la red, cada recorte se normaliza con ecualización adaptativa
de contraste (CLAHE, \textit{Contrast Limited Adaptive Histogram Equalization})
sobre el canal de luminancia del espacio CIELAB, con límite de recorte $2{,}0$ y
retícula de $4 \times 4$ píxeles. El paso no es opcional.

La justificación es empírica y cuantificable. Los recortes reales del parqueadero
son capturas nocturnas: su luminancia mediana es de $76{,}6$ contra $106{,}4$ en
las sintéticas, y su desviación estándar de $20{,}7$ contra $34{,}4$. Sin
normalización, el modelo recibe una distribución de entrada que no vio durante el
entrenamiento. Medido sobre $32$ placas reales con la misma configuración de
entrenamiento y el mismo número de iteraciones, el modelo pasa de $64{,}3\%$ a
$99{,}6\%$ de exactitud por carácter al incorporar CLAHE. En la corrida completa
de dos etapas, la diferencia es de $14\%$ a $67{,}6\%$ de exactitud por carácter
sobre el conjunto de prueba congelado.

El mismo paso se aplica en producción
(\texttt{preprocess()} en \texttt{backend/adapters/outbound/ml/plate\_crnn.py}).
Un reentrenamiento del modelo con CLAHE servido sin CLAHE reintroduce exactamente
el problema que motivó su incorporación, razón por la cual el preprocesamiento
está duplicado en el adaptador y cubierto por una prueba automática.

\subsubsection{Estrategia de transferencia en dos etapas (Two-Stage Transfer Learning)}

Para superar la limitación del tamaño reducido de las muestras reales, el
entrenamiento se ejecuta en dos etapas diferenciadas:

\begin{itemize}
    \item \textbf{Etapa 1 (preentrenamiento sintético):} $8$ épocas sobre
        $359.992$ matrículas sintéticas con formato boliviano. Se congelan las
        capas iniciales del extractor (desde \texttt{conv1} hasta
        \texttt{layer2}) y se habilitan \texttt{layer3} y \texttt{layer4}, con
        tasas de aprendizaje diferenciadas: $3 \times 10^{-4}$ para el backbone
        y $3 \times 10^{-3}$ para la cabeza BiLSTM-clasificador.
    \item \textbf{Etapa 2 (ajuste fino real):} $500$ épocas sobre las $395$
        placas reales, con el backbone completamente descongelado, tasa de
        aprendizaje $1 \times 10^{-3}$ constante y aumento de datos geométrico
        (\textit{crop jitter}): cada recorte se desplaza y escala aleatoriamente
        entre $0{,}90$ y $1{,}12$ antes de normalizarse.
\end{itemize}

Cada una de estas decisiones reemplaza una configuración anterior que fue medida
y descartada, y la sustitución se documenta porque el resultado sin ellas era
sistemáticamente peor:

\begin{itemize}
    \item \textbf{Congelar todo el backbone} en la etapa 1 limita la exactitud por
        carácter al $55\%$: solo $10{,}5$ de $21{,}7$ millones de parámetros
        resultan entrenables, y la cabeza no puede reconfigurar características
        genéricas de ImageNet en representaciones de glifos.
    \item \textbf{Uniformar la tasa de aprendizaje} en $3 \times 10^{-3}$ para
        todos los parámetros provoca el colapso de CTC hacia la secuencia vacía:
        pérdida plana en $2{,}83$ y $0\%$ de coincidencia exacta durante $14$
        épocas consecutivas.
    \item \textbf{Un planificador de tasa por meseta} (\textit{ReduceLROnPlateau})
        en la etapa 2 reduce la tasa a $6 \times 10^{-7}$ en aproximadamente $60$
        épocas, porque la pérdida es ruidosa. El entrenamiento se congela sin
        emitir ningún aviso.
    \item \textbf{Entrenar $30$ épocas reales} equivale a $360$ pasos de
        gradiente para $21{,}7$ millones de parámetros: insuficiente por dos
        órdenes de magnitud frente a la etapa 1.
\end{itemize}

El \textit{crop jitter} merece una nota aparte porque su efecto fue el mayor de
todos. Los recortes de entrenamiento fueron producidos por la caja de PaddleOCR,
ajustada a los glifos efectivamente reconocidos; en producción los recorta el
detector propio, con otra geometría. Esa discrepancia se midió sobre $88$ recortes
de vehículo con texto verificado por un operador: el sistema completo pasó de
$59{,}4\%$ a $73{,}7\%$ de exactitud por carácter al incorporar el aumento
geométrico, y de $9{,}1\%$ a $26{,}1\%$ de coincidencia exacta. Es decir, una
parte sustancial del error atribuido al reconocimiento era en realidad
sensibilidad al encuadre.

\subsubsection{Post-procesamiento sintáctico: medido y desactivado}

El sistema incluye un post-procesamiento determinista
(\texttt{post\_process\_bolivian\_plate}) basado en la gramática oficial, que
reescribe las posiciones numéricas y alfabéticas según las confusiones visuales
más frecuentes.

Su evaluación sobre el conjunto de prueba congelado arroja un resultado contrario
al esperado: reduce la exactitud por carácter de $78{,}9\%$ a $78{,}6\%$. La causa
es que el formato boliviano admite letras dentro del bloque numérico ---por
ejemplo \texttt{216OZAC}--- y la regla convierte una lectura correcta en
incorrecta. En consecuencia, el post-procesamiento queda desactivado por defecto
tanto en el entrenamiento como en el adaptador de producción, y se conserva en el
código únicamente como referencia de la regla sintáctica.

\subsection{Determinación de eventos y trayectorias}

La conversión de trayectorias en eventos de acceso se realiza en el módulo
\texttt{tracking.py}. Se definen dos zonas clave sobre la escena: \texttt{street\_zone}
(vía pública) e \texttt{inside\_zone} (rampa interior del parqueadero), divididas por
la línea virtual de control (\texttt{tripwire}).

El sistema monitorea la posición del punto de contacto con el suelo (centro inferior de
la caja del vehículo). Un vehículo se clasifica como evento de \textit{ingreso} únicamente
si su trayectoria nace o atraviesa la zona de vía y su punto inferior cruza la línea
virtual penetrando en la zona interior. Las trayectorias que permanecen paralelas a la
vía o no alcanzan la zona interior se clasifican como \textit{tránsito} y se guardan como
negativos duros para evitar falsos registros.

\subsection{Clasificación de atributos del vehículo}

Los atributos del vehículo (tipo de carrocería y color) se obtienen procesando los
recortes generados por la detección de YOLO con dos clasificadores ResNet18
independientes, entrenados sobre el mismo pool congelado de atributos y con
aumento de datos diferenciado por tarea. El clasificador de tipo distingue cinco
categorías (automóvil, camioneta, minibús, moto y vagoneta) y admite variación de
color, porque su criterio es la silueta; el de color distingue once categorías y
aplica únicamente variación de exposición, nunca de saturación ni de tono, porque
el tono es la etiqueta que debe predecir. Los resultados medidos de ambos se
reportan en las Subsecciones~3.7.3 y~3.7.4.
